\section*{Capítulo \ref{ch:g12:past-climates} --- Ler os climas do passado}

\begin{solution}{exo:g12:past-climates:1}
(a) um \omterm{def:g12:past-climates:proxy}{indicador indireto}; (b) um \omterm{def:g12:past-climates:ice-core}{testemunho de gelo}; (c) um \omterm{def:g12:past-climates:glacial-cycle}{interglacial};
(d) a \omterm{def:g12:past-climates:pollen}{análise polínica}; (e) o \omterm{def:g12:past-climates:lgm}{Último Máximo Glacial}.
\end{solution}

\begin{solution}{exo:g12:past-climates:2}
(a) O gelo polar: o gás carbônico e o metano do ar da época.
(b) Os sedimentos de lagos ou a turfa: a vegetação, aqui árvores que
precisam de verões quentes. (c) Os sedimentos oceânicos: o volume de gelo
sobre os continentes e a temperatura do \omterm{def:g10:layered-ocean:thermocline}{oceano profundo}. (d) As árvores: as
condições de crescimento, ano a ano.
\end{solution}

\begin{solution}{exo:g12:past-climates:3}
(a) $-50 - (-120) = \qty{70}{m}$ em \num{7000} anos: $\qty{70000}{mm}/\num{7000}
= \qty{10}{mm}$ por ano. (b) $10/3.7 \approx 2.7$: a subida média da
deglaciação foi cerca de 2.7 vezes mais rápida que a de 2006--2018.
\end{solution}

\begin{solution}{exo:g12:past-climates:4}
Um século são dez décadas: $\qty{2.0}{\degreeCelsius}$ por século. Mil anos
são cem décadas: $\qty{20}{\degreeCelsius}$ por mil anos.
\end{solution}

\begin{solution}{exo:g12:past-climates:5}
(a) \num{23000} anos: a precessão do eixo; \num{42000} anos: a inclinação
do eixo; cerca de \num{100000} anos: a forma da órbita.
(b) $\num{1000000}/\num{41000} \approx 24$ ciclos; $\num{1000000}/\num{100000}
= 10$ ciclos.
\end{solution}

\begin{solution}{exo:g12:past-climates:6}
(a) $266.6 + 7.83 \times (-9.4) \approx 266.6 - 73.6 = \qty{193}{ppm}$: só
\qty{3}{ppm} acima dos \qty{190}{ppm} medidos. (b) $266.6 + 7.83 \times 5.5
\approx \qty{310}{ppm}$, \qty{11}{ppm} acima do valor mais alto medido: o
ponto mais quente fica abaixo da reta. (c) $\Delta T = (427 - 266.6)/7.83 \approx
\qty{20}{\degreeCelsius}$. A reta só descreve como as duas grandezas
variaram juntas, em ciclos naturais, dentro da faixa dos dados
(de \qty{-11}{\degreeCelsius} a \qty{+6}{\degreeCelsius}, de 174 a
\qty{299}{ppm}); ela não é uma lei, e um valor muito fora dessa faixa não
pode ser lido nela.
\end{solution}

\begin{solution}{exo:g12:past-climates:7}
(a) As ervas, \qty{84}{\percent}: uma paisagem aberta quase sem árvores,
fria. (b) O carvalho (\qty{42}{\percent}) e a aveleira (\qty{34}{\percent}):
uma floresta de folhosas, com verões quentes. (c) A bétula se espalha
primeiro, depois o pinheiro; a aveleira e o carvalho chegam por último.
(d) As árvores só se espalham tão depressa quanto as suas sementes viajam e
as suas árvores jovens crescem, a partir dos lugares onde sobreviveram ao
frio; um solo também precisa se formar. A floresta pode se atrasar em
relação ao clima.
\end{solution}

\begin{solution}{exo:g12:past-climates:8}
(a) \qty{300}{cm}, ou seja, \qty{3000}{mm}, em \num{8000} anos: $3000/8000 =
\qty{0.375}{mm}$ por ano. (b) $\num{6000} + (470 - 300)/300 \times \num{8000}
\approx \num{6000} + \num{4530} \approx \num{10500}$ anos. (c) $10/0.375
\approx 27$ anos.
\end{solution}

\begin{solution}{exo:g12:past-climates:9}
(a) $7.16 - 4.11 = 3.05$ milhões de quilômetros quadrados, e $3.05/4.11 \approx
\qty{74}{\percent}$ de terras a mais. (b) Sim: no \omterm{def:g12:past-climates:lgm}{Último Máximo Glacial}, o
mar estava 125 a \qty{134}{m} abaixo do de hoje, muito mais que os
\qty{36}{m} que ligam Bornéu à Ásia no mapa. (c) O mapa usa o fundo do mar de hoje, que não
continuou exatamente o mesmo desde então (sedimentos se depositaram sobre
ele, e o terreno se moveu um pouco); a grade também é grosseira, um ponto a
cada 0.1 grau, e o nível mais baixo do mar só é conhecido com uma margem de
cerca de dez metros.
\end{solution}

\begin{solution}{exo:g12:past-climates:10}
(a) O valor do \omterm{def:g12:past-climates:lgm}{Último Máximo Glacial}, \qty{4.92}{\permil}, é o mais alto.
É o esperado: os grandes mantos de gelo da época armazenavam água pobre em
oxigênio-18, o que deixava o oceano, e as conchas que cresciam nele, mais
ricos em oxigênio-18; a água profunda também era mais fria. (b) Para que as
épocas frias e cobertas de gelo fiquem embaixo e as épocas quentes em cima,
como numa curva de temperatura.
\end{solution}

\begin{solution}{exo:g12:past-climates:11}
(a) Uma correlação mostra que os dois registros sobem e descem juntos, perto
de uma reta; ela não mostra qual causou qual, nem que um deles causou o
outro. (b) Um clima mais quente poderia ter liberado gás carbônico no ar; ou
uma terceira causa, como as mudanças da órbita da Terra, poderia ter movido
os dois; ou cada um poderia ter empurrado o outro. (c) A cronologia das
mudanças, qual veio primeiro, a partir de uma datação precisa do gelo e do
seu ar; e uma compreensão de como o gás carbônico age sobre a radiação da
Terra (\cref{ch:g12:climate-change}).
\end{solution}

\begin{solution}{exo:g12:past-climates:12}
(a) $\num{1300000}/\num{41000} \approx 32$ ciclos entre 2.6 e 1.3
milhões de anos atrás; $\num{700000}/\num{98500} \approx 7$ ciclos nos últimos
\num{700000} anos. (b) \num{41000} anos: a inclinação do eixo; cerca de
\num{100000} anos: a forma da órbita. (c) Um ciclo de \num{41000} anos é
amostrado cerca de $\num{41000}/\num{2000} \approx 20$ vezes, o bastante para
desenhá-lo; uma mudança que dura um século cai entre duas amostras, ou dentro
de uma, e quase não deixa vestígio.
\end{solution}

\begin{solution}{exo:g12:past-climates:13}
(a) $192/10 \approx 19$ e $198/10 \approx 20$: cerca de 20 vezes mais rápido.
(b) Contra \qty{4}{ppm} por século: de $192/4 = 48$ a $198/4 \approx 50$
vezes; contra \qty{42}{ppm}: de $192/42 \approx 4.6$ a $198/42 \approx 4.7$
vezes. O aumento de hoje é cerca de 5 a 50 vezes mais rápido que o do PETM.
(c) O aumento é de $2000 - 900 = \qty{1100}{ppm}$: $1100/42 \approx 26$ séculos,
cerca de \num{2600} anos, à taxa mais rápida do PETM; $1100/195 \approx 5.6$
séculos, menos de 600 anos, à taxa de hoje.
\end{solution}

\begin{solution}{exo:g12:past-climates:14}
(a) Dez vezes 1 a \qty{1.5}{\degreeCelsius}: de 10 a \qty{15}{\degreeCelsius}
por mil anos, ou seja, de 0.10 a \qty{0.15}{\degreeCelsius} por década.
(b) A taxa global de hoje, \qty{0.20}{\degreeCelsius} por década, é de 1.3 a
2 vezes o aquecimento mais rápido sobre a Groenlândia. (c) Um único lugar
pode mudar muito mais que a média sobre o globo, em que as mudanças das
diferentes regiões se compensam em parte: uma taxa global se compara com
uma taxa global.
\end{solution}

\begin{solution}{exo:g12:past-climates:15}
(a) \qty{14}{m} a \qty{12}{m} por mil anos levam $14/12 \approx
1.17$ mil anos, cerca de \num{1170} anos: há cerca de \num{9800} anos.
(b) Não: o mar precisava ficar pelo menos \qty{36}{m} abaixo do de hoje
para ligar Bornéu à Ásia; uma vez passado o nível de \qty{-36}{m}, a ponte
de terra ficou submersa, e Bornéu voltou a ser uma ilha. (c) $36/12 = 3$
mil anos, cerca de \num{3000} anos; na verdade, a subida desacelerou depois
de cerca de \num{8200} anos atrás, então levou mais tempo.
\end{solution}

\begin{solution}{pb:g12:past-climates:1}
\textbf{1.} O próprio gelo, pelos seus isótopos pesados, registra a
temperatura; o ar das suas bolhas registra o gás carbônico.

\textbf{2.} $5.5 - (-10.6) = \qty{16.1}{\degreeCelsius}$.

\textbf{3.} $299 - 174 = \qty{125}{ppm}$.

\textbf{4.} Os pontos ficam perto de uma reta: os dois registros subiram e
desceram juntos. Isso não diz qual causou qual.

\textbf{5.} Os registros do passado são desenhados com as épocas mais
antigas de um lado ou do outro; lido no sentido errado, um aquecimento
seria tomado por um resfriamento.

\textbf{6.} Um \omterm{def:g12:past-climates:glacial-cycle}{período glacial} é um intervalo frio dentro de uma glaciação,
em que os mantos de gelo crescem muito além do seu tamanho atual; um
\omterm{def:g12:past-climates:glacial-cycle}{interglacial} é um intervalo quente entre dois \omterm{def:g12:past-climates:glacial-cycle}{períodos glaciais}, com mantos
de gelo próximos dos de hoje ou menores.

\textbf{7.} Num \omterm{def:g12:past-climates:glacial-cycle}{interglacial}, que começou há cerca de \num{11000} anos.

\textbf{8.} $\num{800000}/\num{98000} \approx 8$: cerca de oito ciclos.

\textbf{9.} A forma da órbita.

\textbf{10.} Um ritmo de \num{41000} anos, o período da inclinação do
eixo.

\textbf{11.} $-0.4 - (-9.4) = \qty{9.0}{\degreeCelsius}$.

\textbf{12.} $264 - 190 = \qty{74}{ppm}$.

\textbf{13.} \qty{6}{\degreeCelsius}; a variação antártica, \qty{9.0}{\degreeCelsius},
é maior, cerca de 1.5 vez a global.

\textbf{14.} $\num{18000} - \num{11000} = \num{7000}$ anos.

\textbf{15.} $6/7 \approx \qty{0.86}{\degreeCelsius}$ por mil anos.

\textbf{16.} Mil anos são cem décadas: $0.86/100 \approx
\qty{0.0086}{\degreeCelsius}$ por década.

\textbf{17.} $0.20 \times 100 = \qty{20}{\degreeCelsius}$ por mil anos.

\textbf{18.} $20/0.86 \approx 23$ vezes.

\textbf{19.} $20/1.5 \approx 13$ e $20/1 = 20$: de 13 a 20 vezes.

\textbf{20.} O aquecimento global de hoje é \textbf{cerca de vinte vezes mais
rápido} que o aquecimento ao sair da última glaciação: 23 vezes a sua taxa
média, de 13 a 20 vezes a sua taxa mais rápida.
\end{solution}
